% !TeX root = ../main.tex
\section{Domain Switching and Authentication}
\label{sec:design:impers}
\keypoint{
	\thename's API design allows the introduction of flexible compartmentalization  boundaries to existing software through in-place annotations.
}
However, such a design must be robust against  \emph{domain impersonation} (\srimpers).
% \keypoint{
% In this section, we show how \thename prevents \emph{domain impersonation} (\srimper) with its design.
% }
In \thename's threat model, domain impersonation happens when a compromised domain diverts the control flow into \thename APIs, \eg, calling the domain entry gates with arbitrary arguments, or corrupts \thename's critical values, \eg, \gls{mte} tag ID stored in memory to gain illegal access to domain-private resources.
\thename prevents such domain impersonation using two techniques.
First, \thename \emph{authenticates} the domain entry sites by requiring them to construct \emph{entry tokens}, which the reference monitor can verify before granting the domain switch.
Second, it introduces an in-userspace \emph{authentication} protocol that securely fetches the currently executing domain's \gls{mte} tag.
In both cases, after successful authentication, \thename securely places its critical values (the domain's memory tag and the per-instance modifier) in compiler-reserved registers and fetches them when needed to prevent illegal modifications.
% By forbidding the compiler from using the reserved register and enforcing a coarse-grained CFI, \thename ensures that these registers are used only for their designated purposes.

% Second, \thename uses two reserved registers to store the \gls{mte} tag and the per-instance modifier, preventing the attacker from tampering with its tag and modifier management.

% \hdr{Domain authentication}
% \added{
% Domain authentication enables trustworthy runtime domain ID fetching, which allows us to avoid hardcoding \gls{mte} tag values into functions.
% \thename authenticates the currently executing domain with a data structure called the \gls{dst}.
% Additionally, a global variable that stores the domain ID of the currently active domain (\texttt{curr\_dom}) is also maintained.
% During initialization, \thename's runtime library populates the table with the  \emph{domain signatures} of each declared domain ({\tt DST[i] = }\pac{DBi}{0}{0}).
% \thename then makes the \gls{dst} read-only throughout program execution to prevent tampering.
% }


% \hdr{Authenticated domain ID retrieval}
% \thename enables trustworthy runtime domain ID fetching, which allows us to avoid hardcoding \gls{mte} tag values into functions.
% Hence, \thename-protected functions that employ \gls{mte} tagging can be invoked from \emph{multiple} domains.
% % This authenticated runtime retrieval of domain \gls{mte} tag allows a \thename-protected functions to be called from \emph{multiple} domains.
% % % that have domain-private object authentication (i.e., \texttt{pacdb})
% Domain ID retrieval is performed with the \gls{dst} that is filled with 
% The variable, called \texttt{int curr\_dom}, is updated on each \texttt{capac\_enter(domID)} and \texttt{capac\_exit()} invocation.  

% The following code snippet shown above illustrates how \thename fetches the signature of the current domain and authenticates it. 
% \begin{lstlisting}[style=sslab-capacity-c]
%     domain_signature_t dom_sig = DST[curr_dom];
%     AUTDB(dom_sig); %     // dom_sig now contains /*!\textcolor{clr-comment}{\tagnum{Dom}}!*/ if authentication succeed
%     assert(dom_sig == curr_dom);
%     // /*!\textcolor{clr-comment}{\tagnum{Dom}}!*/ is known from this point onward
% \end{lstlisting}


% The above procedure can be inlined to verify the current domain ID, such that it can be used as a \gls{mte} tag for the pending operation. 
% Notably, \texttt{capac\_malloc()} uses the procedure to authenticate the callee's domain ID and tag the allocated memory accordingly. 
% The procedure is also imperative in runtime domain stack and pointer tagging in \thename's instrumentation, as we will explain in \cref{sec:design:instrumentation}. 



\hdr{Secure domain switching.}
\label{sec:design:libcapac}
\label{sec:design:switch-auth} % \hdr{Bootstraping}
The program enters a domain through the API \texttt{capac\_enter} (\autoref{tab:apis}) that internally invokes an \texttt{ioctl} syscall to the reference monitor.
The reference monitor then switches the currently activating \key{DB} to that of the target domain by writing into the dedicated \gls{pa} key registers with its kernel privilege. Hence, an attacker who can invoke \texttt{ioctl} calls with arbitrary values could achieve \emph{impersonation} and access domain-private references.
\keypoint{
	\thename protects the domain switches from impersonation by authenticating domain entry points with \key{G}.
}
Before a domain switching request, the call gate signs the domain ID argument in \texttt{capac\_enter(domID, mod)} with a \texttt{pacga} instruction, then passes the signed argument to \thename reference monitor.
Since there is no dedicated authentication instruction that uses \key{G}, the reference monitor uses \texttt{pacga} to generate the valid argument with the domain ID and compare it with the argument from the call gate. Domain entry is granted if the two values match.
After the entry is granted, the call gate loads the per-instance modifier into a reserved floating-point register (\texttt{ModReg}), which is to be used by our instrumentation for pointer authentication demonstrated in \cref{sec:instrumentation}.
Finally, the call gate clears the token-containing register to prevent the leaking of the entry token.
The implementation of \texttt{capac\_enter} is shown in \autoref{capacity:entry}, which is to be inlined to the domain switching sites.

\begin{listing}[h]
	\begin{minted}{c}
    PACGA(token, target_id, 0)
    if(!ioctl(capac_fd, CAPAC_ENTER, token, modifier)) 
        exit();
    // Load per-instance modifier into ModReg register
    asm volatile("fmov ModReg, %[mod]" ::[mod] "r"(modifier) :);  
    // Clear entry token from register
\end{minted}
	\caption{\thename domain entry gate}\label{capacity:entry}

\end{listing}

\keypoint{
	\thename endows authenticity to the entry points by preventing attackers from diverting the control flow into arbitrary instructions with \gls{cfi} and ensuring that the instruction \texttt{pacga} is absent within the process but the domain entry sites.
}
\key{G} is currently unused in both deployed and proposed \gls{pa}-based defenses.
Therefore, it is unlikely that it would naturally appear under normal circumstances.
Also, since ARM is a RISC architecture, we can generally rule out the issue of unaligned, unintended occurrences of the instruction, an issue that makes preventing illegal instruction occurrences challenging on x86 architectures \cite{erim}.

% After, the \texttt{capac\_exit} call simply instructs the reference monitor to revert the domain key to that of the ambient domain.

% \added{
% The reference monitor also receives and activates an instance-specific modifier, which it uses as the modifier to sign domain-private file descriptors (\cref{sec:design:ref-mon})  
% }


% \begin{figure}[t]
% 	\centering

% \caption{\thename domain entry call gate}
% \label{fig:dom_callgate}
% \end{figure}
\textbf{Authenticated domain ID retrieval.}
\thename enables a trustworthy runtime domain ID fetching procedure.
This avoids hardcoding \gls{mte} tags into functions that require memory tagging since these functions might be used in confused deputy attacks to tag arbitrary memory regions.
Moreover, fetching the domain ID at runtime allows a \thename-protected function to be called from multiple domains.
The procedure is also imperative in runtime domain memory and pointer tagging in \thename's instrumentation, as we will explain in \cref{sec:design:instrumentation}.
\thename's private heap memory allocator, \texttt{capac\_malloc()}, also uses the procedure to authenticate the callee's domain ID and tag the allocated memory.
% that have domain-private object authentication (i.e., \texttt{pacdb})

To prepare for runtime domain ID retrieval, \thename first constructs a \gls{dst} that is filled with \emph{domain signatures} of each declared domain ({\tt DST[i] = }\pac{DBi}{0}{0}).
After initialization, the table remains read-only throughout program execution to prevent tampering.
Additionally, a global variable that stores the domain ID of the currently active domain, \texttt{int curr\_dom}, is updated on each \texttt{capac\_enter(domID)} and \texttt{capac\_exit()} invocation.

The following snippet is inlined to verify the current domain ID, such that it can be used as a \gls{mte} tag for the pending operation:
% The variable, 
% \begin{listing}[h!]
\begin{minted}[escapeinside=||]{c}
    domain_signature_t dom_sig = DST[curr_dom];
    AUTDB(dom_sig);
    assert(dom_sig == curr_dom);
    // Form a tag mask shifting left by 56 bits
    asm volatile("lsl %[tag], %[tag], 56" : [tag] "=r"(target_id)::); 
    // Load tag mask into TagReg register
    asm volatile("fmov TagReg, %[tag]" ::[tag] "r"(target_id) :);  
    // |\tagnum{Dom}| is known from this point onward
\end{minted}
% \caption{\thename domain ID authentication}
% \end{listing}
The above procedure first fetches the current domain's signature from the \gls{dst}, then authenticates it with the domain key.
If authentication succeeds, a \emph{tag mask} is constructed, \eg \texttt{0x00ff00..00}, which can be applied to an untagged pointer with a bitwise \texttt{OR}.
The tag mask is stored in a reserved floating-point register (\texttt{TagReg}), so the memory tagging operations can efficiently retrieve it.
% \begin{lstlisting}[style=sslab-capacity-c,language={C},keepspaces,numbers=none,frame=none]
% int dom_id = dom_auth(DST[curr_dom]);
% if (dom_id == FAILED) exit();
% // dom_id 
% \end{lstlisting}
% The code snippet shown above illustrates how \thename fetches the signature of the current domain and authenticates it.

% \hdr{Optimization with unused registers}
% After the domain switch is granted, the call gate builds the tag mask \eg, \texttt{0x00ff00..00} which can later be applied on an untagged pointer with a bitwise \texttt{OR}.

% Similarly, the per-instance modifier is stored in another floating-point register.
% Domain private pointer authentication operations later retrieve the modifier from the register to use as the signing/authentication context.




% \added{
% \thename domain switch performs the following actions. 
% Internally, it invokes a \texttt{ioctl} call to the reference monitor to request changing the current domain key to the domain key of \texttt{domID} and allow access to \thename references of that domain.
% Hence, an attacker who can invoke \texttt{ioctl} calls with arbitrary values could achieve \emph{impersonation} and access domain-private references. 
% }

% It signs the target domain ID with \texttt{pacga} as a token of proof that the call originates from a legitimate site, as explained in \cref{sec:design:domains}, and passes the token as an argument to the \texttt{ioctl} call.



% \begin{lstlisting}[style=sslab-capacity-c]
%     PACGA(token, target_id, 0)
%     if(!ioctl(capac_fd, CAPAC_ENTER, token)) 
%         exit();
%     curr_dom = target_id;
% \end{lstlisting}


% \begin{lstlisting}
%    pacga x8, x0, x8      
% \end{lstlisting}




% Enabling trustworthy runtime domain ID fetching allows us to avoid hardcoding \gls{mte} tag values into functions.

% Given our key switching mechanism, \thename must prevents attackers from

% Each domain is authenticated with a per-domain \gls{pa} key called \emph{domain authentication key}, or domain key for short.









% \chcomment{why we need cryptographic}

% \hdr{Domain-private and ambient objects}








% \hdr{Life-cycle protection of program objects}
% The novelty of \thename's design compared to the previous capability systems~\cite{capsicum,cheri-risc,cheri} is that it coherently applies the access control principles throughout the life-cycle of in-process resources.




% A summary of \thename's encapsulated reference types, and their authentication scheme is displayed in \autoref{tbl:objects}\hj{Move}.


% \hdr{Domain memory isolation}
% For domain memory isolation, \thename requires that both the \ptr references and their pointed-to memory region are protected, for that \gls{pa} alone cannot guarantee memory safety. 
% Hence, \thename employs \gls{mte} memory tagging to establish \emph{domain-private} memory. \thename facilities (explained in \cref{sec:implementation}) strictly control the creation and usage of \ptr that could access the tagged memory.
% The ambient domain is given an ID of 0, which means that untagged memory and {\ptr}s are implicitly ambient objects. 
% As shown in \autoref{fig:overview}, the authentication of {\ptr}s utilizes hardware-assisted checks from both \gls{pa} and \gls{mte}.

% \ptr in \thename's security model requires special considerations against pointer forging attacks that reuse \gls{pa} signing instructions~\cite{pac-it-up,kernel-pac}.
% \thename's design incorporates existing \gls{pa}-based \gls{cfi} forward-edge/backward-edge schemes, which is also a common requirement for \gls{pa}-based security~\cite{pac-it-up}.
% Hence, the adversary cannot jump into arbitrary instruction sequences to forge a pointer, and the attack surface is vastly diminished.
% We further discuss the security of \gls{cfi} and the remaining issues (\eg, \emph{signing oracles}) in \cref{sec:discussion:oracle}.


% \hdr{Reference delegation} \thename also supports references \emph{delegation}, \ie, the transferrtherring of resources access right from one domain into another.
% This is an essential feature to support multi-domain interaction in complex programs.
% Our webserver example discussed in \cref{sec:evaluation:porting-app} illustrates the necessity of such domain-to-domain reference delegation; a domain that interacts with the private server key for TLS handshake produces a session key, which s then delegated to the session handling domain.





% We will further explain the design and implementation of the components that enforce complete mediation in the next section.

% \hdr{Security requirements}




